\documentclass[11pt,a4paper]{article}

\usepackage[margin=2.3cm]{geometry}
\usepackage{amsmath,amssymb}
\usepackage{booktabs}
\usepackage{graphicx}
\usepackage{xcolor}
\usepackage{hyperref}
\usepackage{listings}
\usepackage{enumitem}
\usepackage{array}

\hypersetup{
  colorlinks=true,
  linkcolor=blue!55!black,
  urlcolor=blue!55!black
}

\lstset{
  basicstyle=\ttfamily\small,
  frame=single,
  breaklines=true,
  columns=fullflexible,
  backgroundcolor=\color{gray!5}
}

\newcommand{\code}[1]{\texttt{\detokenize{#1}}}
\newcommand{\vect}[1]{\boldsymbol{#1}}
\newcommand{\Normal}{\mathcal{N}}

\title{Topology-Informed Multimodal MPPI for SLIP\\
\large Sandbox and Nav2 Implementation Progress: V0--V8}
\author{MSc thesis development notes}
\date{\today}

\begin{document}
\maketitle
\tableofcontents
\newpage

\section{Purpose of this document}

This document records what has been implemented in the Python/NumPy sandbox,
the independent PyTorch port, and the ROS~2 Nav2 C++ integration. It explains
why each version was introduced, how the pieces interact, and what remains
before deployment on the physical SLIP robot.

The present contribution is not merely a geodesic gradient that shifts one
Gaussian. The intended contribution is:

\begin{quote}
A topology-informed multimodal proposal generator for MPPI that maps distinct
local free-space branches into dynamically feasible, geometry-adaptive control
distributions.
\end{quote}

The sandbox implements V0 through V7.2, including static and dynamic
benchmarks and an independent PyTorch execution path. V6 remains an
analytical, configurable SLIP prototype rather than a calibrated hardware
model. The Nav2 C++ port has reached V8 Phase~3: the controller constructs a
finite geodesic body, extracts up to three pseudopods, converts them into
time-varying ancillary control sequences, and injects promise-weighted sample
groups into Nav2 MPPI. Physical parameter identification, collision validation
of the ancillary rollouts, asynchronous field reconstruction, and disciplined
Gazebo/robot benchmarks remain open.

\section{Robot and deployment context}

\begin{table}[h]
\centering
\begin{tabular}{>{\bfseries}p{0.34\linewidth}p{0.56\linewidth}}
\toprule
Platform & SLIP, approximately 70\,kg skid-steer mobile robot \\
Software & ROS~2 Humble, Ubuntu 22.04, Nav2 \\
Compute & NVIDIA Jetson Orin Nano \\
Localization & RoboSense Airy 3D LiDAR with FAST-LIO2 \\
Command interface & Differential-drive $(v,\omega)$ velocity command \\
Physical behavior & Skid steer, therefore turning and slip differ from an
ideal differential-drive model \\
Typical MPPI configuration & $T=56$, $\Delta t=0.05$\,s, batch size near 2000 \\
Environment & Narrow, cluttered, scaffold and construction-like spaces \\
\bottomrule
\end{tabular}
\caption{Target deployment context.}
\end{table}

\section{System overview}

The implemented data flow is:

\begin{center}
\fbox{\parbox{0.91\linewidth}{
\centering
costmap/obstacles $\rightarrow$ local geodesic body
$\rightarrow$ pseudopod branches (V1)
$\rightarrow$ branch control means (V2)
$\rightarrow$ grouped mixture proposals (V3)
$\rightarrow$ corrected mixture weights (V4)
$\rightarrow$ ESS guard (V4.1)
$\rightarrow$ geometry-adaptive covariance (V5)
$\rightarrow$ footprint/skid-steer feasibility (V6)
$\rightarrow$ committed MPPI mode
$\rightarrow [v,\omega]$

\medskip
The global A* path supplies a protected fallback proposal and normal-navigation
reference. It is not replaced by the pseudopods.
}}
\end{center}

Three controllers are publicly selectable:

\begin{table}[h]
\centering
\begin{tabular}{p{0.18\linewidth}p{0.72\linewidth}}
\toprule
\textbf{Name} & \textbf{Role} \\
\midrule
\code{vanilla} & Standard single-distribution MPPI baseline. \\
\code{local} & Local geodesic-body guidance with optional pseudopod extraction
and grouped sampling. \\
\code{hybrid} & Global-path tracking plus gated local assistance. This is the
closest sandbox analogue of the intended Nav2 behavior. \\
\bottomrule
\end{tabular}
\end{table}

Earlier ray, full-map flow, and waypoint-commit variants were removed from the
runnable API. Their useful lessons remain in the local/hybrid design, while
their old figures are preserved under \code{artifacts/images/legacy}.

\section{MPPI baseline}

Let a control sequence be
\begin{equation}
  U = [u_0,\ldots,u_{T-1}], \qquad u_t = [v_t,\omega_t]^\top.
\end{equation}
For the sandbox prediction model,
\begin{align}
  \theta_{t+1} &= \theta_t + \omega_t\Delta t,\\
  x_{t+1} &= x_t + v_t\cos(\theta_{t+1})\Delta t,\\
  y_{t+1} &= y_t + v_t\sin(\theta_{t+1})\Delta t.
\end{align}

The trajectory cost contains collision, footprint-clearance repulsion, and
goal/path/geodesic terms:
\begin{equation}
  S(U) = S_{\mathrm{collision}} + S_{\mathrm{clearance}}
       + S_{\mathrm{navigation}}.
\end{equation}
Vanilla MPPI forms weights
\begin{equation}
  w_k =
  \frac{\exp[-(S_k-S_{\min})/\lambda]}
       {\sum_j\exp[-(S_j-S_{\min})/\lambda]}
\end{equation}
and updates the nominal sequence using the weighted sampled controls.

\section{Version history}

\subsection{V0: local geodesic-field-guided MPPI}

V0 constructs a local ``amoeba body'' with two floods:

\begin{enumerate}[leftmargin=*]
  \item An unweighted, radius-limited flood from the robot defines reachable
  free cells. This is the body's shape.
  \item A viscosity/obstacle-weighted navigation flood inside that body defines
  a cost-to-go field.
\end{enumerate}

The field influences MPPI through a steering bias and a geodesic trajectory
cost. The field is rebuilt every few controller cycles. Reachability and
viscosity are deliberately separate: viscosity is a route preference, not a
sensing-range definition.

V0 is useful guidance, but it still describes one proposal distribution and
therefore does not preserve explicit left/right route alternatives.

\subsection{V1: topology-separated pseudopod extraction}

V1 extracts up to three branches from the local body's actual reachable
frontier:

\begin{enumerate}[leftmargin=*]
  \item Rank promising frontier cells using reach and navigation promise.
  \item Apply spatial suppression to avoid many nearby candidates.
  \item Trace each candidate back toward the robot through the reach field.
  \item Reject geometrically equivalent centrelines.
  \item Retain separate branches only when they genuinely diverge around dry
  or obstructed space.
  \item Match branch identities between field rebuilds using
  \code{BranchTracker}.
\end{enumerate}

The result is a collection of explicit objects $P_0,P_1,P_2$, each with an
endpoint, centreline, score, and stable ID. Open space intentionally collapses
to one geometric branch; a pillar can produce left and right branches.

\subsection{V2: branch-to-control-sequence conversion}

Each centreline is converted into a full $[T,2]$ mean sequence:
\begin{equation}
  \mu_m = [u_{m,0},\ldots,u_{m,T-1}], \qquad
  u_{m,t}=[v_{m,t},\omega_{m,t}]^\top.
\end{equation}

The converter uses receding lookahead, heading error, curvature feed-forward,
clearance-aware speed reduction, and the sandbox's velocity and acceleration
limits. Every proposal is rolled through the same dynamics used by MPPI.
It reports feasibility, minimum clearance, progress, and tracking error.
Infeasible branches are not promoted into active sampling modes.

\subsection{V3: fixed grouped multimodal sampling}

V3 forms the proposal family
\begin{equation}
  q(U) = \sum_{m=1}^{M}\pi_m
  \Normal(U;\mu_m,\Sigma_m).
\end{equation}
At this stage $\Sigma_m$ uses the same fixed diagonal standard deviations for
all modes. A deterministic fraction of the batch is assigned to each feasible
pseudopod and to an always-present fallback. Samples retain their mode labels
through rollout and scoring.

Only the selected mode's controls are averaged for the output update. This is
important: directly averaging good left-route and right-route control
sequences may generate a bad trajectory through the obstacle between them.
V3 uses uncorrected within-mode MPPI weights and is retained as an ablation
baseline.

It is also the weighting rule frozen for the final thesis benchmark. With
$M_k$ active pseudopods, the protected A* fallback receives
$\pi_0=0.25$ of the batch and each pseudopod receives
$\pi_m=0.75/M_k$. If no pseudopod is exposed, the fallback receives the full
batch. For mode $m$, weights are normalized only over samples carrying its
label,
\begin{equation}
  w_{i\mid m}=\frac{\exp[-(S_i-S_{m,\min})/\lambda]}
  {\sum_{j:z_j=m}\exp[-(S_j-S_{m,\min})/\lambda]},
\end{equation}
and only the committed mode $m^*$ updates the output sequence,
\begin{equation}
  U^*=\sum_{i:z_i=m^*}w_{i\mid m^*}U_i.
\end{equation}
Raw mixture-importance correction and adaptive covariance are disabled in the
frozen result configuration because their ablations did not establish a
reliable aggregate advantage.

\subsection{V3.1: grouped-sampling stability and allocation}

The first grouped controller required additional policy to remain continuous
over repeated control cycles. V3.1 added:
\begin{itemize}[leftmargin=*]
  \item a protected allocation for the global-path or nominal fallback;
  \item a persistent shifted warm-start sequence for every active mode;
  \item a free-energy estimate and effective sample size (ESS);
  \item mode selection with minimum dwell time, an improvement margin, and
  multi-cycle confirmation;
  \item significant angular sign-flip, mode-switch, fallback-selection,
  path-error, safety, and timing diagnostics.
\end{itemize}

The fallback allocation is configured by \code{--fallback-share}. The selected
development default reserves 25\% of the batch for the protected path/fallback
and divides the remaining 75\% among feasible branches. Mode commitment is
configured by \code{--mode-dwell}, \code{--mode-confirm}, and
\code{--switch-margin}.

\subsubsection{V3.1 development ablation}

The ablation used BARN worlds 4, 48, and 49 with seeds 0 and 1, robot radius
0.34\,m, $K=1024$, $T=56$, and $\Delta t=0.05$\,s. It compared the original
single-Gaussian hybrid baseline against four grouped policies:
\begin{description}[leftmargin=2.7cm,style=nextline]
  \item[Path25] 25\% protected path allocation.
  \item[Path50] 50\% protected path allocation.
  \item[Path70] 70\% protected path allocation.
  \item[Conservative] 50\% path allocation, 20-cycle dwell, five-cycle
  confirmation, and a 0.6 switching margin.
\end{description}

\begin{table}[h]
\centering
\resizebox{\textwidth}{!}{%
\begin{tabular}{lrrrrrrrrr}
\toprule
\textbf{Policy} & \textbf{Success} & \textbf{Time} & \textbf{Path} &
\textbf{Worst clear.} & \textbf{Flips} & \textbf{Switches} &
\textbf{Fallback} & \textbf{ESS} & \textbf{Path error} \\
\midrule
Baseline & 6/6 & 24.5\,s & 9.28\,m & 0.098\,m & 166.7 & 0.0 &
1.00 & --- & 0.102\,m \\
Path25 & 6/6 & 20.3\,s & 8.98\,m & 0.095\,m & 15.8 & 8.3 &
0.68 & 17.3 & 0.064\,m \\
Path50 & 6/6 & 20.3\,s & 9.00\,m & 0.091\,m & 15.5 & 6.7 &
0.73 & 19.3 & 0.066\,m \\
Path70 & 6/6 & 20.4\,s & 8.99\,m & 0.091\,m & 17.3 & 6.7 &
0.78 & 19.8 & 0.065\,m \\
Conservative & 6/6 & 20.4\,s & 8.99\,m & 0.083\,m & 17.5 & 3.2 &
0.53 & 18.7 & 0.072\,m \\
\bottomrule
\end{tabular}}
\caption{V3.1 grouped-sampling development ablation. ESS is the mean
within-mode effective sample size; fallback is the selected-cycle fraction.}
\label{tab:v31ablation}
\end{table}

Path25 was selected as the development default. Relative to the baseline, it
reduced significant angular sign reversals by about 90\%, reduced mean path
error by about 37\%, and completed the six cases roughly four seconds faster.
Increasing the protected path allocation did not monotonically improve
tracking or clearance. The conservative policy reduced switching but had the
worst clearance and path error, showing that excessive commitment can preserve
a degrading mode.

These results were not treated as proof of superiority. There were only six
runs per configuration, grouped worst-case clearance remained slightly below
the baseline, all worlds came from the same environment family, and no dynamic,
stale-plan, or skid-steer experiment was included.

\subsubsection{V3.1 observations that shaped later versions}

\begin{enumerate}[leftmargin=*]
  \item Minimum per-mode ESS approached one in every grouped setting. This
  directly motivated log-domain V4 weighting, explicit degeneracy tests, and
  the later V4.1 ESS guard.
  \item Episodes showed only 2--12 committed switches but 6--12 distinct branch
  IDs. Endpoint-based identity matching was not stable enough for a long-lived
  Nav2 proposal bank.
  \item Mean grouped field construction was about 27\,ms, branch proposal
  generation 6.4\,ms, and path proposal generation 4.2\,ms on the development
  workstation. This motivated the asynchronous topology-worker design for the
  Nav2 port.
  \item The circular footprint and ideal differential-drive model remained
  insufficient representations of the SLIP skid-steer platform.
\end{enumerate}

\subsection{V3.2: global replanning and gated assistance}

The hybrid controller has two operational states:
\begin{description}[leftmargin=2.8cm,style=nextline]
  \item[NORMAL] Follow the protected global-path proposal while the path is
  locally valid and the robot is progressing.
  \item[ASSIST] Expose pseudopod modes when the path is blocked, clearance is
  confined, or progress stalls.
\end{description}

The sandbox can periodically rerun A* from the current pose. A successful new
plan updates the same controller instance, increments a plan version, and
returns the system to NORMAL when appropriate. This models the fact that Nav2's
global planner replans as the robot moves; the controller is not permanently
fixated on the original path.

\subsection{V4: corrected Gaussian-mixture importance weighting}

V4 implements the requested correction:
\begin{equation}
  \log \widetilde{w}_k =
  -\frac{S(U_k)}{\lambda}
  + \log p(U_k)
  - \log q(U_k).
  \label{eq:v4weight}
\end{equation}

\paragraph{Target distribution $p$.}
The nominal target is a diagonal Gaussian around the protected fallback mean:
\begin{equation}
  p(U)=\Normal(U;\mu_{\mathrm{fallback}},\Sigma).
\end{equation}
For \code{hybrid}, the fallback is the dynamically feasible global-path
proposal. For \code{local}, it is the shifted nominal MPPI sequence.

\paragraph{Proposal distribution $q$.}
The complete density includes every active branch and fallback:
\begin{equation}
 q(U)=\sum_{m=1}^{M}\pi_m
 \Normal(U;\mu_m,\Sigma).
\end{equation}
The coefficients $\pi_m=N_m/K$ equal the actual fixed-group sample
proportions. The log mixture is evaluated stably:
\begin{equation}
 \log q(U)=\operatorname{LSE}_{m}
 \left[\log\pi_m+\log\Normal(U;\mu_m,\Sigma)\right],
\end{equation}
where $\operatorname{LSE}(a_i)=a_{\max}+
\log\sum_i\exp(a_i-a_{\max})$.

\paragraph{Latent versus projected controls.}
Gaussian densities are evaluated on raw latent sequences before velocity and
acceleration projection. Projection is deterministic but many-to-one; treating
the clipped/projected controls as ordinary Gaussian samples would give an
incorrect density.

\paragraph{Preserving route modes.}
The corrected weights determine mode evidence and the conditional update
inside each mode. The committed-mode policy remains active, so the estimator
does not average topologically incompatible routes into one command sequence.

\paragraph{Important caution.}
Mathematical correction does not guarantee lower finite-sample variance.
The $\log p-\log q$ ratio can reduce effective sample size when branch means
are far from the target. V3 and V4 must therefore be compared empirically;
V4 should not be claimed superior solely because it is formally corrected.

\subsection{V4.1: cycle-coherent ESS guard}

The blocked-path experiment confirmed the variance concern. Some branch groups
contained hundreds of samples but had corrected ESS near one. The resulting
mode evidence favored the protected fallback and raw V4 failed to pass the
blockage.

V4.1 therefore evaluates the corrected effective sample size for each active
mode $m$:
\begin{equation}
  \operatorname{ESS}_m =
  \frac{1}{\sum_{k \in m} \bar{w}_{k,m}^{\,2}},
\end{equation}
where $\bar{w}_{k,m}$ are the corrected weights normalized within mode. A mode
triggers the guard when
\begin{equation}
  \operatorname{ESS}_m <
  \max\left(8,\;0.02N_m\right).
  \label{eq:essguard}
\end{equation}

If any active mode triggers Equation~\ref{eq:essguard}, the complete control
cycle uses V3 within-mode weights and V3 free energies. Falling back for the
whole cycle is intentional: using V3 evidence for one mode and corrected V4
evidence for another would compare quantities on different statistical
scales. Raw V4 weights, free energies, and ESS values are retained as
diagnostics even during a guarded cycle.

\subsubsection{Development comparison}

The first normal-path comparison on BARN world 48, seed 0, was not a meaningful
mixture test. Automatic gating remained in NORMAL, so all 1024 samples belonged
to the path component. In that single-component case $p=q$, and V3 and V4
correctly produced identical trajectories.

A genuine multimodal comparison used world 48, seed 0, radius 0.34\,m, a
deliberately blocked centre-lane path, automatic assistance, Nav2-like path
validity, and no replanning:

\begin{table}[h]
\centering
\begin{tabular}{lrrrrrrr}
\toprule
\textbf{Weights} & \textbf{Result} & \textbf{Time} & \textbf{Path} &
\textbf{Clear.} & \textbf{Stall} & \textbf{Switch} & \textbf{Flips} \\
\midrule
V3 & success & 29.00\,s & 10.43\,m & 0.100\,m & 0.60\,s & 6 & 42 \\
V4 & timeout & 60.00\,s & 7.56\,m & 0.181\,m & 17.20\,s & 1 & 31 \\
V4.1 & success & 26.95\,s & 10.06\,m & 0.110\,m & 0.95\,s & 6 & 27 \\
\bottomrule
\end{tabular}
\caption{Single-case blocked-path development comparison.}
\label{tab:v41blocked}
\end{table}

V4.1 restored success, but the guard was active for 498 cycles and 1298 of
1355 mode evaluations (95.8\%). This demonstrates a successful degeneracy
fallback, not an advantage from corrected weighting. V3 remains the controller
baseline; raw V4 remains diagnostic-only until a multi-world, multi-seed
benchmark supports a stronger conclusion.

\section{Global planner relationship}

The pseudopod layer is not a replacement global planner. The intended
responsibility split is:

\begin{itemize}[leftmargin=*]
  \item The global planner provides long-range route intent and replans as the
  robot moves.
  \item Nav2 path handling transforms and prunes the latest plan.
  \item MPPI normally follows a dynamically feasible path proposal.
  \item Local pseudopods become additional proposals when blockage,
  confinement, or local-minimum behavior makes the normal proposal unreliable.
  \item If topology data is stale or invalid, the controller immediately falls
  back to ordinary Nav2 MPPI/path behavior.
\end{itemize}

This makes the research claim narrower and more defensible: local topology
provides explicit alternative proposal modes where a single local control
distribution is insufficient.

\section{Software architecture}

\begin{table}[h]
\centering
\small
\begin{tabular}{p{0.28\linewidth}p{0.63\linewidth}}
\toprule
\textbf{File} & \textbf{Responsibility} \\
\midrule
\code{mppi.py} & MPPI rollout, costs, V3 grouped update, V4 correction, mode
selection and warm-start bank. \\
\code{amoeba_controllers.py} & Local and hybrid controller policy, field
rebuild, path gating and proposal preparation. \\
\code{env.py} & BARN environment, clearance, global and local flow fields. \\
\code{pseudopods.py} & V1 frontier selection, centreline tracing,
deduplication and stable IDs. \\
\code{proposals.py} & V2 centreline-to-control conversion and feasibility. \\
\code{grouped_sampling.py} & Group allocation, sampling, constraint projection,
Gaussian log densities, mixture log density, V4 importance weights. \\
\code{robot_model.py} & Shared V6 ideal/skid motion, oriented footprint, and
turn/speed safety margins. \\
\code{astar.py} & Sandbox global A* planner. \\
\code{simulation.py} & Episode loop and performance metrics. \\
\code{cli.py} & Command-line configuration and renderer selection. \\
\code{visualization.py} & Matplotlib static figures and compatibility animation. \\
\code{pyqtgraph_visualization.py} & Fast PyQtGraph live animation. \\
\code{tests/test_importance_weighting.py} & Closed-form and integration tests for V4. \\
\code{experiments/v5_ablation.py} & Fixed/adaptive covariance benchmark harness. \\
\code{experiments/v6_ablation.py} & Ideal/analytical-SLIP matched benchmark harness. \\
\code{torch_port/} & Independent batched tensor core, BARN/geodesic runtime,
A* hybrid controller, tests, benchmark, and PyQtGraph observer. \\
\bottomrule
\end{tabular}
\caption{Current sandbox code map.}
\end{table}

\section{Verification status}

At the time of writing, 54 automated tests pass in the PyTorch environment.
The dependency-light system-Python workflow passes the original tests and
skips nine optional PyTorch tests. They cover:

\begin{itemize}[leftmargin=*]
  \item two branches around a pillar and one branch in open space;
  \item branch-ID stability after small robot motion;
  \item full bounded V2 sequences and matching rollout dynamics;
  \item collision-based proposal infeasibility;
  \item deterministic sample allocation and preserved mode labels;
  \item velocity, acceleration and angular-acceleration projection;
  \item persistent per-mode warm starts and switching hysteresis;
  \item global-path fallback protection and NORMAL/ASSIST transitions;
  \item closed-form diagonal Gaussian log densities;
  \item manually evaluated two-component mixture density;
  \item numerical stability in a far Gaussian tail;
  \item reduction to vanilla MPPI weights when $p=q$;
  \item V4 controller diagnostics and normalized weights;
  \item V4.1 ESS-triggering, cycle-coherent fallback, and diagnostics.
  \item V5 component-specific sampling, mixture densities, covariance
  schedules, and ablation bookkeeping.
  \item V6 ideal/slip integration, footprint clearance, safety margins, and
  ablation bookkeeping;
  \item deterministic NumPy/PyTorch rollout, projection, cost, density, and
  importance-weight parity;
  \item independent PyTorch BARN, A*, NORMAL/ASSIST, protected-path, and
  blocked-path assistance behavior.
\end{itemize}

The tests demonstrate implementation consistency, not hardware readiness.
Single-world development runs are not sufficient evidence of robustness.

\section{How to run the sandbox}

\subsection{Run all tests}
\begin{lstlisting}
cd /home/saran/robohouse_ws/amoeba_sandbox
python3 -m unittest discover -v
\end{lstlisting}

\subsection{Recommended V4 visualization}
\begin{lstlisting}
python3 run.py --controller hybrid --world 48 --radius 0.34 \
  --grouped-sampling --importance-weighting \
  --gate auto --nav2-path-validity \
  --replan-every 1.0 --mode-debug-every 20 --animate
\end{lstlisting}

\subsection{Compare V3, V4, and V4.1}

Run the same case three times. V3 omits \code{--importance-weighting}; V4 includes
it, and V4.1 additionally includes \code{--importance-ess-guard}. Use the same
world, robot radius, and seed. Disable animation when measuring wall-clock
computation.

\begin{lstlisting}
# V3
python3 run.py --controller hybrid --world 48 --seed 0 \
  --radius 0.34 --grouped-sampling --gate auto \
  --nav2-path-validity --replan-every 1.0

# V4
python3 run.py --controller hybrid --world 48 --seed 0 \
  --radius 0.34 --grouped-sampling --importance-weighting \
  --gate auto --nav2-path-validity --replan-every 1.0

# V4.1
python3 run.py --controller hybrid --world 48 --seed 0 \
  --radius 0.34 --grouped-sampling --importance-weighting \
  --importance-ess-guard --gate auto --nav2-path-validity \
  --replan-every 1.0
\end{lstlisting}

To exercise a real mixture rather than the single path component, replace the
A* plan with the blocked development path by adding \code{--lane center} and
omitting \code{--replan-every}.

Images belong under \code{artifacts/images}; numerical experiment tables belong
under \code{artifacts/results}. Reusing an output filename overwrites it.

\section{Known limitations and unclear assumptions}

\begin{enumerate}[leftmargin=*]
  \item V6 contains an analytical skid-steer response and rectangular
  footprint, but its dimensions, yaw gain, speed-dependent yaw loss, and
  uncertainty margins have not been identified from the physical SLIP robot.
  \item V3 and raw V4 use fixed diagonal covariance. V5 adapts diagonal
  entries using predicted clearance and turning demand, but the mapping is
  hand-designed and not calibrated from SLIP data.
  \item Mode-conditioned commitment is a policy decision layered on top of the
  corrected importance weights. It intentionally avoids route averaging, but
  its effect must be included in ablations.
  \item The sandbox path cost is only an approximation of Nav2 critics. The C++
  port should reuse PathAlign, PathFollow, obstacle, footprint and constraint
  behavior rather than duplicate them.
  \item Local field reconstruction is synchronous in Python. It must not block
  the Nav2 control loop.
  \item PyQtGraph is the default live renderer; Matplotlib remains available
  for static publication figures and compatibility animation. Both are
  callback observers and are irrelevant to controller timing. Headless runs
  remain authoritative.
  \item V4 has exhibited high importance-weight variance in the blocked-path
  test. V4.1 preserves safety behavior by reverting to V3, but its 95.8\%
  fallback rate means corrected weighting is not yet an active controller
  improvement. ESS and tail behavior require multi-world, multi-seed
  measurement.
\end{enumerate}

\section{V5, V6, and next versions}

\subsection{V5: geometry-adaptive diagonal covariance}

For each mode and horizon step, construct
\begin{equation}
  \Sigma_{m,t} =
  \operatorname{diag}(\sigma^2_{v,m,t},\sigma^2_{\omega,m,t}).
\end{equation}
V5 evaluates clearance along each mode's predicted mean rollout and combines
it with normalized turning demand. Open straight motion widens linear
exploration. Confinement reduces linear variance but retains or increases
angular exploration, since contracting both channels prevented lateral
detours in the first development rule. Sampling and V4 density evaluation use
the same component-specific covariance schedules.

The fixed/adaptive ablation contains 48 matched runs: three worlds, two seeds,
four static/dynamic and normal/blocked scenarios, and two covariance policies.

\begin{table}[h]
\centering
\resizebox{\textwidth}{!}{%
\begin{tabular}{llrrrrrrrr}
\toprule
\textbf{Scenario} & \textbf{Cov.} & \textbf{Success} & \textbf{Coll.} &
\textbf{Time} & \textbf{Worst clear.} & \textbf{Stall} & \textbf{Flips} &
\textbf{Switches} & \textbf{ESS} \\
\midrule
Normal & fixed & 6/6 & 0 & 20.4 & 0.085 & 0.0 & 9.5 & 0.0 & 89.1 \\
Normal & V5 & 6/6 & 0 & 20.4 & 0.095 & 0.0 & 8.8 & 0.0 & 94.7 \\
Blocked & fixed & 6/6 & 0 & 25.8 & 0.076 & 0.1 & 12.7 & 7.5 & 60.0 \\
Blocked & V5 & 6/6 & 0 & 28.3 & 0.091 & 0.7 & 12.3 & 10.8 & 70.3 \\
Dynamic & fixed & 5/6 & 1 & 24.4 & -0.007 & 0.5 & 12.0 & 5.5 & 45.9 \\
Dynamic & V5 & 5/6 & 1 & 23.7 & -0.006 & 0.8 & 12.3 & 5.0 & 60.4 \\
Dyn. blocked & fixed & 4/6 & 2 & 29.3 & -0.007 & 0.8 & 10.5 & 10.3 & 40.5 \\
Dyn. blocked & V5 & 4/6 & 1 & 33.5 & -0.011 & 3.6 & 13.0 & 15.0 & 71.0 \\
\bottomrule
\end{tabular}}
\caption{V5 fixed/adaptive covariance ablation. Times are means over
successful runs. The V5 dynamic-blocked failures are one collision and one
timeout; fixed covariance has two collisions.}
\label{tab:v5ablation}
\end{table}

V5 improves worst static clearance and mean ESS in every scenario, but does
not improve aggregate success. Blocked cases are slower and exhibit more mode
switching. The two policies also fail on different dynamic cases. V5 is
therefore complete as a research implementation and ablation, but remains an
experimental option; fixed covariance with V3 weights remains the deployment
baseline because the analytical V6 model has not improved aggregate success.
Full covariance is deliberately deferred.

\subsection{V6: analytical SLIP-specific feasibility}

V6 centralizes prediction, proposal rollout, episode execution, path validity,
and collision metrics in one configurable robot model. For commanded
$u_t=[v_t,\omega_t]^\top$, the provisional skid-steer yaw response is
\begin{equation}
  \omega_t^{\mathrm{eff}}=
  \frac{k_\omega\omega_t}
       {1+k_s|v_t|/v_{\max}},
  \qquad k_\omega=0.82,\quad k_s=0.55,
\end{equation}
followed by
\begin{align}
  \theta_{t+1}&=\theta_t+\omega_t^{\mathrm{eff}}\Delta t,\\
  x_{t+1}&=x_t+v_t\cos(\theta_{t+1})\Delta t,\\
  y_{t+1}&=y_t+v_t\sin(\theta_{t+1})\Delta t.
\end{align}
The assumed footprint is a $0.90\,\mathrm{m}\times0.65\,\mathrm{m}$ oriented
rectangle sampled along its boundary and centre. Signed footprint clearance is
\begin{equation}
  c_f(x_t)=\min_{p\in\mathcal F(x_t)}c_{\mathrm{map}}(p),
\end{equation}
and the hard analytical uncertainty allowance is
\begin{equation}
  m_t^{\mathrm{extra}}=
  0.08\frac{|\omega_t^{\mathrm{eff}}|}{\omega_{\max}}
  +0.03\frac{|v_t|}{v_{\max}}.
\end{equation}
A rollout is marked colliding when $c_f(x_t)<m_t^{\mathrm{extra}}$. The
orientation-free A* and geodesic layers use the width proxy
\begin{equation}
  r_p=\frac{0.65}{2}+0.08=0.405\ \mathrm{m};
\end{equation}
the full circumcircle made all three development worlds unreachable, while
oriented rollout checks still enforce the complete rectangle.

The matched 48-run comparison used worlds 4, 48, and 49; seeds 0 and 1; and
normal, blocked, dynamic, and dynamic-blocked scenarios.
\begin{table}[h]
\centering
\begin{tabular}{llrrr}
\toprule
\textbf{Scenario} & \textbf{Model} & \textbf{Success} & \textbf{Coll.} &
\textbf{Worst clear.} \\
\midrule
Normal & ideal & 6/6 & 0 & 0.085 \\
Normal & analytical SLIP & 5/6 & 1 & $-0.000$ \\
Blocked & ideal & 6/6 & 0 & 0.076 \\
Blocked & analytical SLIP & 5/6 & 1 & $-0.019$ \\
Dynamic & ideal & 5/6 & 1 & $-0.007$ \\
Dynamic & analytical SLIP & 4/6 & 2 & $-0.013$ \\
Dynamic blocked & ideal & 4/6 & 2 & $-0.009$ \\
Dynamic blocked & analytical SLIP & 2/6 & 4 & $-0.027$ \\
\bottomrule
\end{tabular}
\caption{V6 ideal/analytical-SLIP development ablation. Clearance is in
metres and uses each model's own footprint convention.}
\label{tab:v6ablation}
\end{table}

V6 successfully establishes shared model and footprint plumbing and exposes
contacts hidden by the circular ideal model. It does not improve success and
is not hardware-ready. The dynamic failures motivate V6.1 horizon-aware
obstacle prediction; physical bags and footprint measurements are required for
V6.2 calibration.

\subsection{Independent PyTorch acceleration track}

The sibling \code{torch_port} implementation leaves the NumPy reference
unchanged. CPU topology operations construct the BARN map, A* plan, rolling
geodesic body, pseudopods, and ancillary means. Device tensors perform grouped
sampling, constraint projection, batched rollout, footprint lookup, trajectory
costs, and MPPI updates. Its hybrid layer protects the A* path in NORMAL,
activates pseudopods in ASSIST, applies dwell/confirmation commitment, and can
replan periodically. The light-mode PyQtGraph renderer is an observer and is
excluded from timing. CUDA throughput and multi-world episode batching must be
reported separately from mathematical controller comparisons.

\subsection{V6.1 and V6.2}

V6.1 replaces current-position-only dynamic collision scoring with a
horizon-indexed obstacle estimate. For obstacle $j$ with observed centre
$o_{j,0}$ and velocity estimate $\hat v_j$,
\begin{equation}
  \hat o_{j,t}=o_{j,0}+\hat v_j t\Delta t,
\end{equation}
and rollout clearance becomes a function of both sample and horizon index.
The implemented predictor holds its last extrapolation after 2.0,s and grows
the obstacle radius by 0.03,m per predicted second. It is opt-in, so static
behavior and the V6 baseline remain unchanged. A first world-42, seed-0,
15,s A/B run timed out in both cases, but prediction increased minimum
clearance from 0.050,m to 0.246,m and progress from 6.55,m to 6.77,m. This
is promising evidence, not a success claim; multi-world V7 statistics are
still required. A subsequent two-world paired pilot (worlds 42 and 48,
one seed, 256 samples, 300 control cycles) reduced collisions from one to
zero, raised mean minimum clearance from 0.072,m to 0.237,m, and added only
about 0.2,Hz of measured overhead. Both predicted runs timed out, so the
pilot supports safety but does not yet establish task success. V6.2 then
identifies footprint,
actuation, slip, acceleration, and stopping parameters from physical SLIP
measurements rather than further hand-tuning the analytical placeholders.

\subsection{V7: benchmark and ablation}

The benchmark infrastructure compares vanilla MPPI, path-biased MPPI, and the
complete amoeba controller. V7.1 adds dynamic-junction cases aimed directly at
the thesis claim, and V7.2 adds the PyTorch execution track. Existing results
are development evidence rather than a final thesis comparison: the final
matrix must use matched worlds and seeds and report success, collision rate,
time, path length, minimum clearance, stall duration, angular sign flips, mode
switches, fallback fraction, ESS, field/proposal/MPPI timing, and memory use.

\subsection{Evaluation metrics and reporting protocol}

The thesis evaluation will not treat every terminal value as an independent
claim.  Metrics are grouped by the question they answer and are evaluated on
matched world--seed pairs.  Let episode $e$ contain robot position $p_t$,
commanded linear velocity $v_t$, control period $\Delta t$, termination time
$T_e$, and shortest collision-free A* reference length $L_e^{*}$.  The primary
task metric is
\begin{equation}
 s_e=\mathbb I[\text{goal reached without collision before the time limit}],
 \qquad \hat S=\frac{1}{E}\sum_{e=1}^{E}s_e .
 \label{eq:benchmark-success}
\end{equation}
Failures must also be separated into static collision, moving-obstacle
collision, robot--robot collision, and timeout; success rate alone can conceal
whether a controller is safe but excessively conservative.

Safety is reported using signed footprint clearance, for which zero denotes
contact and a negative value denotes overlap.  For obstacle class
$q\in\{\mathrm{static},\mathrm{moving}\}$,
\begin{equation}
 c^{\min}_{q,e}=\min_{0\leq t\leq T_e}c_q(x_t),
 \qquad
 c^{\min}_{\mathrm{peer},e}
 =\min_t\bigl(\lVert p_t^1-p_t^2\rVert_2-r_1-r_2\bigr).
 \label{eq:benchmark-clearance}
\end{equation}
The peer expression is the circular clearance proxy currently used by the
reciprocal runner; static and scripted-obstacle checks use the implemented
SLIP footprint geometry.  Clearance distributions and their lower tail are
more informative than reporting only their mean.

Efficiency and recovery behaviour are measured by
\begin{align}
 L_e &= \sum_{t=1}^{T_e}\lVert p_t-p_{t-1}\rVert_2,
 &\rho_{L,e}&=L_e/L_e^{*},\\
 L_e^{-} &= \sum_t\max(0,-v_t)\Delta t,
 &T_{\mathrm{stall},e}&=\Delta t\sum_t
 \mathbb I[|v_t|<v_{\epsilon}\wedge d_{\mathrm{goal}}>d_{\epsilon}].
 \label{eq:benchmark-efficiency}
\end{align}
Completion time, path length, and reverse distance are already logged.  The
normalised path ratio and stall duration are derivable or require a small
telemetry addition.  Mode switches, topology switches, recovery activations,
and planned-yield duration diagnose oscillation and deadlock rather than task
quality directly; lower values are desirable only when success and clearance
are maintained.

For topology-aware experiments the runner also records the number of route
candidates, chosen route indices, predicted route-pair clearance, feasibility,
active topology, route-conflict prior, prediction horizon, and topology-switch
count.  These measurements test the causal claim that an alternative route was
available, recognised, and used.  An early-decision latency---time from episode
start until a conflicting route is rejected---should be added before the final
multi-robot study.

MPPI proposal health is characterised by normalised effective sample size
\begin{equation}
 \mathrm{ESS}=\frac{(\sum_{k=1}^{N}w_k)^2}{\sum_{k=1}^{N}w_k^2},
 \qquad \mathrm{nESS}=\mathrm{ESS}/N,
 \label{eq:benchmark-ess}
\end{equation}
plus fallback fraction, allocation per mode, and selected mode.  These are
mechanism metrics: they explain a result but do not replace success, safety,
or efficiency.  ESS and fallback telemetry exist in the NumPy research path
but must be exported consistently by the PyTorch benchmark before being used
in the final comparison.

Under peer-intention prediction, the mean hypothesis probabilities and their
entropy are reported,
\begin{equation}
 H(p)=-\sum_h p_h\log p_h .
 \label{eq:benchmark-entropy}
\end{equation}
High entropy indicates that the predictor remains uncertain; it is not by
itself either good or bad.  Prediction quality should ultimately be evaluated
against the realised peer motion using negative log likelihood or a proper
calibration score, not agreement with the old heuristic.

Computational performance uses mean, 95th-percentile, and maximum controller
latency, together with
\begin{equation}
 R_{\mathrm{sim}}=T_e/T_{\mathrm{wall}},
 \qquad
 f_{\mathrm{miss}}=\frac{1}{T_e}\sum_t
 \mathbb I[t_{\mathrm{controller}}>\Delta t].
 \label{eq:benchmark-runtime}
\end{equation}
The current reciprocal runner logs controller latency and wall time, so
$R_{\mathrm{sim}}$ is derivable; deadline-miss fraction and peak GPU memory
remain additions.  Animation is disabled for timing trials.

The minimum thesis comparison therefore reports, per controller and scenario:
success with its failure breakdown, minimum static/dynamic clearance,
completion time, path length and normalised path ratio, reverse and stall
behaviour, controller latency, and the relevant mechanism metrics.  Continuous
metrics are presented as median and interquartile range with a paired
bootstrap confidence interval; paired success outcomes use McNemar's test.
Worlds and seeds, robot model, sample count, horizon, device, and time limit
remain identical within each controller comparison.  Exploratory pilot runs
are labelled separately and are not pooled into the confirmatory benchmark.

\subsection{Reciprocal interaction-risk ablation}

The PyTorch reciprocal pilot evaluates four peer futures: constant velocity,
following the peer's A* path, yielding to a stop, and reversing. For MPPI
sample $k$, let $J_k^{(h)}$ denote its interaction cost under hypothesis $h$.
Three aggregation rules were tested:
\begin{align}
  J_k^{\mathrm{worst}} &= \max_h J_k^{(h)},\\
  J_k^{\mathrm{mean}} &= \frac{1}{H}\sum_{h=1}^{H}J_k^{(h)},\\
  J_k^{\mathrm{CVaR}_{0.5}} &=
    \operatorname{mean}\!\left(\text{worst half of }\{J_k^{(h)}\}_{h=1}^H\right).
\end{align}
The matched five-seed static and two-moving-obstacle suites used world 93,
the SLIP footprint, 1024 samples, a 56-step horizon, and CUDA. Constant-
velocity prediction succeeded in 10/10 trials. Worst-case hypotheses
succeeded in 6/10 and CVaR in 7/10. Static/dynamic mean minimum peer clearance
was 0.127/0.182,m for constant velocity, 0.680/0.693,m for worst-case, and
0.793/0.706,m for CVaR. The hypothesis variants also reversed farther and
switched modes much more frequently.

Equal-weight mean aggregation was then screened only on the two previously
difficult seeds to avoid an unnecessary full matrix. It succeeded in one of
four static/dynamic trials; both dynamic trials timed out with 38 and 69 mode
switches. Aggregation alone therefore does not solve the deadlock. A single
collision hypothesis retains a contribution near $10000/4=2500$ even under
the mean, while unsupported yield and reverse futures remain active.

The next interaction model will retain hard static and scripted-obstacle
safety checks but weight peer futures using observation-conditioned
probabilities,
\begin{equation}
  J_k = J_{\mathrm{static},k}+J_{\mathrm{path},k}
        +\sum_{h=1}^{H}p_h J_{\mathrm{peer},k}^{(h)}.
\end{equation}
Path alignment, deceleration, and observed backward motion will update the
corresponding probabilities. Reverse will become a hysteretic deadlock-
recovery state rather than an ordinary continuously competing mode. This
hand-crafted probabilistic model is the intermediate ablation before a
learned local-goal probability estimator.

The first hand-crafted stage is implemented: each path or pseudopod ancillary
mean is rolled out as an ego local intention, assigned a separate temporally
smoothed peer-intention distribution, and its MPPI samples use those weights.
The runner records mean intention probabilities and entropy. Difficult-seed
validation and a hysteretic reverse-recovery state remain pending. The first
difficult-static screen passed both seeds: mean completion time was 25.93,s,
minimum peer clearance 0.327,m, mode switches 8.5, and reverse distance
0.48,m. Equal-weight mean had passed only one of the same two seeds, with 23.5
switches and 2.80,m reverse distance. This supports the probability-weighted
architecture as a screening result; the matched dynamic screen remains next.
The subsequent dynamic screen passed seed 0 but seed 1 timed out with 72 mode
switches and 4.475,m of reversing even though inferred reverse probability was
below one percent. This identified the independently competing reverse
proposal as the remaining failure mechanism. Reverse is now gated by a
hysteretic \code{NORMAL/YIELD/RECOVERY\_REVERSE/REJOIN} state machine with
persistent-conflict entry, bounded recovery, and cooldown. Only dynamic seed 1
must be repeated before expanding the matrix.
The forward-priority repeat passed: both robots finished in 49.65,s, the non-
yielding robot used zero reverse motion, and peer clearance was 0.166,m. The
yielding robot nevertheless invoked recovery five times, total mode switches
were 36, and static clearance reached 0.027,m. This is a passed causal
regression, not a final efficiency or robustness result.

\subsection{V7.3: topology-guided parallel Amoeba MPPI}

The first multi-route implementation selected one pair of A* routes before
the episode and discarded the alternatives. That preselection is useful as a
screen, but it is not topology-driven parallel trajectory optimization. V7.3
therefore keeps every accepted global route active as a separate MPPI mode.
This section records the implemented approximation and explicitly does not
claim equivalence to T-MPC: route geometry is currently generated once from
the static map, whereas generation of new dynamic space--time homotopy classes
online remains future work.

For robot $a$, repeated penalized A* produces the candidate set
\begin{equation}
  \mathcal R^a=\{R^a_0,\ldots,R^a_{K_a-1}\}.
\end{equation}
After accepting route $R_j$, cells within a tube of radius $r_{\mathrm{pen}}$
receive an additional traversal cost before the next A* search. Numerical
variants of the same route are rejected using the symmetric mean nearest-point
separation
\begin{equation}
 d_{\mathrm{sym}}(R_i,R_j)
 =\frac{1}{2|R_i|}\sum_{p\in R_i}\min_{q\in R_j}\|p-q\|_2
 +\frac{1}{2|R_j|}\sum_{q\in R_j}\min_{p\in R_i}\|q-p\|_2.
 \label{eq:global-route-separation}
\end{equation}
A candidate is retained only if $d_{\mathrm{sym}}\ge\epsilon_R$ for every
accepted route. This is a practical geometric diversity proxy, not a formal
homotopy signature. A single-corridor world may consequently and correctly
return only one spatial route.

Before closed-loop execution, the reciprocal pilot screens every route pair
and an initial yielding delay $\delta$ for robot~1. With constant-speed
arc-length interpolation $p(R,t;v,\delta)$, it selects
\begin{equation}
 (i_0,j_0,\delta_0)
 =\arg\min_{i,j,\delta}
 \left[
   \frac{L(R^1_i)+L(R^2_j)}{v_R}
   +w_\delta\delta
   +J_{\mathrm{pair}}(R^1_i,R^2_j,\delta)
 \right].
 \label{eq:initial-route-pair}
\end{equation}
This supplies an initial compatible pair only. It does not remove the other
routes from either controller.

For every retained route $R_k$, the same receding branch follower used for a
pseudopod constructs a route-conditioned ancillary mean
\begin{equation}
  \mu_{R_k}=\mathcal C(R_k,x_t,u_{t-1})
  =[u_{k,0},\ldots,u_{k,T-1}].
\end{equation}
The controller's active proposal set is consequently
\begin{equation}
 \mathcal Z_t=
 \underbrace{\{R_0,\ldots,R_{K-1}\}}_{\text{global route modes}}
 \cup
 \underbrace{\mathcal A_t}_{\text{local amoeba pseudopods}}
 \cup
 \underbrace{\mathcal Y_t}_{\text{wait/yield/reverse modes}}
 \cup\{0\},
 \label{eq:parallel-topology-set}
\end{equation}
where mode $0$ is the nominal fallback. Samples are never averaged across
these labels. Instead,
\begin{equation}
 U_n\mid z_n=z\sim\mathcal N(\mu_z,\Sigma),
 \qquad \sum_{z\in\mathcal Z_t}N_z=N.
 \label{eq:parallel-topology-mixture}
\end{equation}

Route samples are scored against their own reference, not against the
currently selected route. Let $d_k(x)$ be lateral distance to $R_k$ and let
$\mathcal L_k(x)$ be remaining route length at the closest projected point.
The implemented route term for sample trajectory $X_n$ is
\begin{equation}
 J_{R_k}(X_n)=
 5\,[d_k(x_{n,T})+\mathcal L_k(x_{n,T})]
 +2\,\overline{[d_k+\mathcal L_k]}(X_n)
 +4\,\overline d_k(X_n).
 \label{eq:route-conditioned-cost}
\end{equation}
Each label then has its own conditional weights, update, and free energy,
\begin{align}
 w_{n\mid z}
 &=\frac{\exp[-(S_n-S_{z,\min})/\lambda]}
         {\sum_{m:z_m=z}\exp[-(S_m-S_{z,\min})/\lambda]},\\
 \widehat U_z&=\sum_{n:z_n=z}w_{n\mid z}U_n,\\
 F_z&=S_{z,\min}
 -\lambda\log\!\left[
   \frac{1}{N_z}\sum_{n:z_n=z}
   \exp\!\left(-\frac{S_n-S_{z,\min}}{\lambda}\right)
 \right].
 \label{eq:parallel-route-free-energy}
\end{align}
The selected topology changes only after the candidate is better by a margin
$\Delta_F$, satisfies a minimum dwell time, and wins for $C$ confirmation
cycles:
\begin{equation}
 z_{t+1}=
 \begin{cases}
 z', & F_{z'}+\Delta_F<F_{z_t}
       \text{ for $C$ cycles and dwell is satisfied},\\
 z_t, & \text{otherwise}.
 \end{cases}
 \label{eq:topology-hysteresis}
\end{equation}
Only $\widehat U_{z_{t+1}}$ supplies the command. This prevents a left-route
mean and right-route mean from producing an unsafe averaged command.

The local MPPI horizon is $T\Delta t=2.8$\,s, which is too short to expose a
head-on conflict near the middle of a long corridor. V7.3 therefore adds a
second, cooperative route timescale. Let $\bar R^a_k(x_t)$ denote the
uncompleted suffix of candidate route $k$, and let $\bar R^b$ be the peer's
announced active route suffix. At nominal speeds $v_a$ and $v_b$,
\begin{align}
 T_{k}^{\mathrm{route}}
 &=\max\!\left(
   \frac{L(\bar R^a_k)}{v_a},
   \frac{L(\bar R^b)}{v_b}\right),\\
 p^a_k(t)&=p(\bar R^a_k,t;v_a,0),
 \qquad
 p^b(t)=p(\bar R^b,t;v_b,0),\\
 d_{k}(t)&=\|p^a_k(t)-p^b(t)\|_2.
\end{align}
With required centre separation
\begin{equation}
 d_{\mathrm{req}}=r_a+r_b+m_R,
 \qquad
 e_k(t)=\max(0,d_{\mathrm{req}}-d_k(t)),
\end{equation}
the implemented long-horizon conflict prior is
\begin{equation}
 J_k^{\mathrm{route\mbox{-}conflict}}
 =w_c\left[
   \max_{0\le t\le T_k^{\mathrm{route}}}e_k(t)^2
   +\frac{1}{4}\,\overline{e_k(t)^2}
 \right],
 \qquad
 (v_a,m_R,w_c)=(0.30,0.12,80).
 \label{eq:long-route-conflict}
\end{equation}
The peak term prevents a short predicted collision from being diluted by a
30-second route, while the mean term penalizes sustained overlap. This scalar
is added to every MPPI sample assigned to route $R_k$ before computing
$F_{R_k}$. Thus eventual route conflict is visible near the beginning of the
episode while short-horizon MPPI retains responsibility for detailed dynamics,
footprint clearance, and control constraints.

The current implementation provides two complementary topology levels:
global route labels $R_k$ and local geodesic pseudopods $P_m$. It does not yet
construct a separate geodesic body for every $R_k$, nor does it generate new
LEFT/RIGHT/BEHIND/AHEAD space--time classes around changing obstacle
predictions. Those are the next steps required before describing the method as
a complete adaptation of topology-driven parallel trajectory optimization.

\subsection{Space-time topology: Phase 0 and first wiring}

A first, independent answer to the LEFT/RIGHT/BEHIND/AHEAD gap above now
exists in the reference NumPy sandbox (\texttt{spacetime.py}), separate from
the PyTorch reciprocal-pilot route machinery described above. It floods free
\emph{space-time} rather than free space: a moving obstacle's predicted path
becomes a tube in $(x,y,t)$ that a route can pass before, wait out, or detour
around, giving those as genuinely distinct candidate modes instead of
something MPPI can only discover via cost-shaped sampling noise around one
existing branch mean. The full construction --- the time-expanded state
space, the wait-cost-variation mechanism used to build two non-degenerate
routes, the synchronized-time distinctness test, and the trigger/local-goal
conditions --- is derived in
\texttt{docs/amoeba\_mathematics\_guide.tex}, Section~``Space-time topology
extension''; it is not repeated here.

Wired into \texttt{mppi.py}'s \texttt{build\_branch\_proposals} as
\texttt{\_spacetime\_alternatives}: when a predicted obstacle crosses a
branch's own centreline continuation, a feasible ``wait'' and/or ``detour''
route is added to that cycle's proposal set with a distinct, stable
\texttt{branch\_id} (offset by $+100$/$+200$), so it is picked up automatically
as an ordinary mode with no changes to mode selection. Off by default;
confirmed to leave the proposal set unchanged even with a genuine predicted
crossing present when disabled.

Two implementation mistakes were found and fixed before this triggered at
all, both instances of a check that runs, looks plausible, and silently
answers the wrong question --- the same failure class as the exact-footprint
and plant/plan-mismatch fixes earlier this cycle. First, an initial
``two-route'' construction bounded which time layer counted as arrival
(an early deadline or a late one); this degenerates whenever reaching the
goal safely is already possible well before the late deadline, since the
cheapest way to satisfy it is then simply to arrive early and idle --
physically the same route, wrongly labelled distinct. Second, the local
search goal was first computed using the controller's own $v_{\max}$, while
the search grid's actual implied speed is $\Delta r_s/\Delta t_\ell$
(generally coarser); a goal picked as reachable at $v_{\max}$ was often not
reachable by the grid that had to reach it, and separately, the crossing
\emph{detector} was first limited to the base proposal's short MPPI-horizon
rollout, so a longer search horizon changed nothing because the detector
never looked that far ahead. Fixed by (i) varying the relative cost of
waiting versus detouring, verified to produce two routes that are genuinely
different in a synthetic crossing scenario (one held position for five
layers with a $0.2\,$m lateral excursion, the other never waited and swung
$0.6\,$m aside), and (ii) checking the crossing against the branch's own
nominal-speed centreline continuation, evaluated out to the full search
horizon, rather than either the achieved rollout or $v_{\max}$.

Verified end to end on a synthetic controller-level scenario (disabled
leaves the proposal set untouched even with a crossing present; enabled with
no crossing adds nothing; a genuine crossing yields exactly the expected
extra modes, consumed correctly by mode selection; a short closed-loop
episode with both space-time modes live completes with finite commands
throughout) and, separately, confirmed to trigger in practice: a headless
scan across several BARN worlds with default oscillating-obstacle settings
found dozens of crossings and added modes over short episodes on most of
them, not just in the hand-tuned synthetic test. Not yet validated against
the V6.1 dynamic-obstacle and V7.1 dynamic-junction scenarios above, which
changes what a mode \emph{is} and is required before any performance claim;
the path fallback and multiple simultaneous obstacles remain unhandled.

\subsection{Final blocked-path ancillary attribution}

The final blocked-split experiment freezes the SLIP footprint, 1024 samples,
a 56-step horizon, fixed covariance, and stable V3 grouped weights. Vanilla
MPPI failed 0/5 through static collision, path-tracking MPPI failed 0/5 through
collision with the scripted blocker, and full Amoeba MPPI succeeded 5/5. To
separate geodesic guidance from pseudopod sampling, a flow-only controller kept
the identical field, path bias, gate, dynamics, and costs while disabling
ancillary proposals.

Across ten paired seeds, flow-only succeeded 2/10, with one peer collision,
two static collisions, and five timeouts. Full Amoeba succeeded 9/10, with no
collisions and one timeout. Seven pairs favored Amoeba, none favored flow-only,
two succeeded under both, and one timed out under both. The exact two-sided
McNemar result is $p=0.015625$. Amoeba reduced mean stall time from 19.96,s to
5.02,s, increased mean peer clearance from 0.131,m to 0.703,m, and selected
pseudopod modes in 47.6\% of cycles. Mean controller time rose from 137.56,ms
to 148.68,ms in the parallel outcome run.

The remaining Amoeba failure was a safe seed-5 timeout with 18.1,s stall time
and 19 mode switches. It is retained without tuning. This result supports the
specific claim that topology-separated ancillary distributions add robustness
beyond a geodesic flow field alone in the declared blocked-path scenario;
ordinary narrow-world generality and isolated serial timing remain separate
experiments.

\subsection{Static generality and the single-path biased baseline}

The ten-world, five-seed static matrix produced 25/50 successes for vanilla
MPPI, 39/50 for path-tracking MPPI, 45/50 for single-path Biased MPPI, and
50/50 for full Amoeba. The respective collision counts were 25, 11, 5, and
zero. Mean clearance was 0.216,m, 0.233,m, 0.251,m, and 0.258,m. Mean
successful completion time was 24.21,s, 25.36,s, 19.86,s, and 20.87,s,
respectively.

This result must be interpreted separately from the blocked-path attribution.
In the static matrix Amoeba had ancillary modes available during 70.9\% of
cycles, but selected a pseudopod mode in 0\% of cycles; ASSIST occupancy was
only 2.0\%. Thus the static result supports the complete gated geodesic/path
hybrid, not a causal pseudopod-sampling claim. The blocked-split experiment is
the mechanism evidence.

The label path MPPI in the original matrix denotes an A* path-tracking cost;
it is not itself Biased MPPI. The \code{path\_biased} comparator instead uses
one A*-derived ancillary control mean and draws the full MPPI batch around that
mean. It improved six matched trials over path tracking without losing any.
Its five failures all occurred in world~7; Amoeba succeeded in those same five
trials and never lost a matched trial. Since those discordances share one
geometry, they show a repeatable world-specific advantage, not five independent
environment-level observations. The exact paired two-sided result is
$p=0.0625$. This completes the comparison between vanilla sampling, path-cost
shaping, single-path biased sampling, and multi-ancillary Amoeba MPPI.

\begin{figure}[htbp]
  \centering
  \includegraphics[width=\textwidth]{artifacts/images/experiments/final/static_success_clearance.png}
  \caption{Final static generality comparison over ten BARN worlds and five
  matched seeds per controller. Negative minimum clearance denotes a
  footprint collision.}
  \label{fig:final-static-generality}
\end{figure}

\begin{figure}[htbp]
  \centering
  \includegraphics[width=0.84\textwidth]{artifacts/images/experiments/final/world7_trajectory_comparison.png}
  \caption{Matched world-7, seed-0 replay. Single-path Biased MPPI follows its
  unique reference into a footprint collision, whereas geodesic Amoeba MPPI
  takes the viable side of the terminal obstacle cluster and succeeds. The
  blue body and flow are displayed only for Amoeba because they are inactive
  in the single-path baseline.}
  \label{fig:world7-trajectory}
\end{figure}

\subsection{Isolated reference-Python timing}

Controller timing was measured sequentially, with no parallel workers, on
critical worlds 7 and 48 and comfortable world 42. All controllers used seed
0, 1024 samples, and a 56-step horizon. Vanilla, path-tracking, single-path
biased, and Amoeba MPPI required mean update times of 31.45,ms, 109.60,ms,
133.80,ms, and 147.97,ms, respectively. Their mean episode-level 95th
percentiles were 35.25,ms, 115.98,ms, 166.20,ms, and 194.70,ms; the worst
Amoeba update was 226.97,ms.

Consequently, the synchronous NumPy reference does not satisfy a 20,Hz
(50,ms) deployment budget. Amoeba adds 14.2,ms, or 10.6\%, over the closer
single-path Biased MPPI baseline; much of the total latency is already incurred
by path processing and the 1024-by-56 rollout batch. The deployment design
therefore computes the geodesic field asynchronously at low rate and uses a
vectorized C++ or GPU rollout path. These Python measurements characterize the
research implementation and are not claimed as Nav2 plugin timing.

\subsection{Offline learned pseudopod ranking}

An optional learned extension retains deterministic collision-free geometry
and predicts only context-dependent utility over feasible pseudopods. A shared
CNN receives robot-centred occupancy, signed clearance, body, membrane,
navigation-value, and candidate-mask channels; an MLP adds proposal and
relative-goal features. Its utility score is made risk-aware using the
auxiliary collision head,
\begin{equation}
  s_m^{\mathrm{safe}}
  =\hat{s}_m-\beta\,\sigma(\hat{c}_m),
  \qquad \beta=20,
\end{equation}
where $\sigma(\hat{c}_m)$ is predicted collision probability. The coefficient
and checkpoint were frozen using validation worlds before blind evaluation.

The dataset contains all 134 SLIP-feasible BARN worlds, split by complete
world into 99 training, 18 validation, and 17 blind-test worlds. On 459 blind
multimodal snapshots, the frozen ranker achieved 387/459 (84.3\%) oracle
top-1 agreement, compared with 358/459 (78.0\%) for geometric promise. Mean
surrogate utility regret decreased from 0.377 to 0.0156. The ranker selected
9 collision-labelled candidates compared with 16 for the heuristic and
avoided all seven cases in which the heuristic selected a colliding candidate
despite the presence of a safe alternative.

This is offline counterfactual ranking evidence, not closed-loop navigation or
robot collision-rate evidence. The checkpoint is not allowed to command or
reallocate MPPI samples. It has now been connected to the PyTorch controller in
observation-only shadow mode: inference runs only during multimodal
\code{ASSIST} cycles and logs the learned branch, geometric-promise branch,
collision probabilities, agreement, and latency. The shadow result is excluded
from sample allocation, MPPI costs, mode selection, and actuation. A
same-seed regression test confirms exact command parity with shadow inference
disabled. The next validation stage measures latency and online agreement on
development worlds, followed---only if that gate passes---by a conservative
blend with protected fallback and unchanged footprint rejection.

The shadow panel passed on four development worlds: all paired outcomes and
complete state traces were identical, while 1,245 multimodal cycles required a
mean of 1.96\,ms for ranker inference. Learned and geometric choices agreed on
63.3\% of those cycles. Consequently, a controlled allocation-only ablation is
enabled. The frozen baseline divides the pseudopod budget equally; the new
geometric policy uses reach-plus-navigation promise, and the conservative
policy blends probabilities as
\begin{equation}
  \pi_m^{\mathrm{blend}}
  =0.75\,\pi_m^{\mathrm{geo}}+0.25\,\pi_m^{\mathrm{learned}}.
\end{equation}
A 10\% uniform mixture floor prevents proposal starvation and 25\% of all
samples remain reserved for the fallback. The learned output changes neither
feasibility checks, trajectory costs, MPPI weighting, nor final mode selection.

Across three seeds and four development worlds, every allocation policy
completed all 12 episodes without collision. The initial one-seed advantage of
the blend did not persist: relative to equal allocation, its mean completion
time increased by 0.363\,s and mean minimum clearance decreased by
4.6\,mm. Fully learned allocation reduced mean time by only 0.175\,s and mode
switches by 0.25 while reducing clearance by 1.5\,mm. These small and
inconsistent paired differences do not support a closed-loop improvement
claim. Learned ranking is therefore retained as an exploratory negative
ablation and observation-only extension; equal allocation remains the primary
Amoeba configuration for the hybrid global-planner benchmark.

Dynamic inspection revealed that periodic global replanning could adopt a
partially executed local detour. Global replanning is therefore suspended while
the hybrid gate is in \code{ASSIST}. In addition, pseudopod $m$ receives
\begin{equation}
 J_m^{\mathrm{rejoin}}
 =w_d\,\overline{d}(\xi_m,\mathcal P)
 +w_r\max\!\left(0,R_m(T)-R_{A^*}(T)\right),
 \qquad (w_d,w_r)=(3,2),
\end{equation}
where $\overline d$ is mean lateral separation from protected path
$\mathcal P$ and $R(T)$ is remaining path length at the rollout endpoint. The
protected A* proposal has zero prior. This discourages locally safe but
globally long pseudopods without disabling local assistance.

The hybrid gate is additionally made predictive. For upcoming path point
$p_i$ at arc length $s_i$, the estimated arrival time is
\begin{equation}
 t_i=\frac{s_i}{\max(|v|,v_{\min}^{\mathrm{gate}})},
 \qquad v_{\min}^{\mathrm{gate}}=0.20\,\mathrm{m/s}.
\end{equation}
Assistance activates before a stall if, for any predicted obstacle $j$,
\begin{equation}
 \|p_i-\hat o_j(t_i)\|
 \le r_{\mathrm{robot}}+r_j+\sigma t_i.
\end{equation}
Thus the safety cost may still command a stop when no bypass is feasible, but
alternative pseudopod sampling begins before all nominal-path samples become
unsafe.

The global planner uses static map clearance; moving objects remain a local
predictive responsibility. After eight-connected A*, line-of-sight
shortcutting, and 5\,cm densification, deterministic Laplacian updates smooth
the polyline. Each update is accepted only when its point and adjacent
midpoints retain footprint clearance. The raw piecewise-linear A* path remains
available as an ablation, and identical smoothing is used for vanilla and
hybrid comparisons.

For transient blockage, the mixture additionally contains an explicit
zero-velocity waiting proposal. MPPI compares waiting with topology-distinct
bypasses rather than stopping only as a side effect of collision penalties.
Let $b_k$ denote the predicted corridor-blocked flag. Dynamic assistance
returns to the protected path after
\begin{equation}
 b_{k-2}=b_{k-1}=b_k=0,
\end{equation}
without requiring displacement during the wait. The three-cycle confirmation
suppresses chatter; the progress anchor and stall counter reset on release.

\subsection{V8: Nav2 C++ integration}

The C++ integration extends \code{nav2_mppi_controller}; it does not replace
Nav2 path handling, critics, motion models, lifecycle management, or command
constraints. Phase~1 introduced shadow-mode field construction and debug
markers. Phase~2 enabled geodesic biasing while retaining a vanilla fallback.
Phase~3 adds up to three pseudopod-derived, time-varying ancillary control
sequences. With the current batch size $N=2000$, the ancillary fraction is
$\alpha=0.20$, so at most 400 samples are divided among the active pseudopods
and approximately 1600 samples retain ordinary Nav2 MPPI sampling. If
$h_m$ is pseudopod $m$'s promise value, its normalized allocation is
\begin{equation}
  \pi_m = \frac{\exp(-h_m/\tau_p)}
                 {\sum_j \exp(-h_j/\tau_p)},
  \qquad \tau_p=0.75,
\end{equation}
followed by integer rounding that preserves the total ancillary budget.
The current operating configuration uses
\code{amoeba_shadow_mode=false}, \code{amoeba_bias_enabled=true},
\code{amoeba_bias_strength=0.2}, and keeps the separate flow critic disabled.

The field is still reconstructed synchronously. A later phase should move
costmap/path processing to a 2--5\,Hz worker that publishes a timestamped,
immutable snapshot for non-blocking consumption by the controller loop.

\section{Nav2 port sketch}

\begin{center}
\fbox{\parbox{0.88\linewidth}{
costmap and transformed/pruned global path\\
$\downarrow$\\
topology/ancillary construction (currently synchronous)\\
$\downarrow$\\
promise-weighted ancillary sample ranges plus vanilla fallback\\
$\downarrow$\\
Nav2 MPPI noise generation and motion prediction (15--20\,Hz)\\
$\downarrow$\\
existing Nav2 critics and footprint collision checks\\
$\downarrow$\\
existing MPPI update, constraints and output filtering
}}
\end{center}

The current port provides sample ranges, pseudopod promises, ancillary paths,
and debug visualization. The next implementation step is to collision-check
the actual ancillary rollouts with Nav2's footprint/costmap machinery and
publish those rollouts separately from their geometric reference centrelines.
Stable mode IDs, per-mode warm starts, free energy/ESS telemetry, timestamps,
and plan/costmap versioning remain desirable for the later asynchronous port.

The visualization interface deliberately separates data by purpose:
\code{/amoeba_debug} is a \code{MarkerArray} for the membrane and field,
\code{/amoeba/ancillary_path_1} through \code{_3} are \code{nav_msgs/Path}
reference centrelines, and \code{/trajectories} is a throttled batched
\code{LINE_LIST} for MPPI samples. The ancillary topics require an RViz Path
display, not a MarkerArray display.

\section{Progress checklist}

\begin{tabular}{p{0.10\linewidth}p{0.68\linewidth}p{0.14\linewidth}}
\toprule
\textbf{Version} & \textbf{Deliverable} & \textbf{Status} \\
\midrule
V0 & Single local geodesic-field-guided MPPI baseline & Complete \\
V1 & Extract and visualize 2--3 topology-separated branches & Complete \\
V2 & Convert each branch to a feasible full control mean & Complete \\
V3 & Fixed grouped multimodal sampling & Complete \\
V3.1 & Allocation, warm starts, commitment, and stability ablation & Complete \\
V3.2 & Gated assistance and periodic global replanning & Complete \\
V4 & Stable full-mixture importance correction & Complete \\
V4.1 & Cycle-coherent ESS degeneracy guard & Complete \\
V5 & Geometry-adaptive diagonal covariance and ablation & Complete \\
V6 & Analytical SLIP skid-steer and footprint feasibility & Prototype \\
V6.1 & Horizon-aware dynamic-obstacle prediction & Prototype complete \\
V6.2 & Physical SLIP parameter identification & Pending \\
V7 & Vanilla/path/Amoeba and dynamic-junction benchmark framework & Prototype complete \\
V7.2 & Independent PyTorch execution and visualization port & Prototype complete \\
V7.3 & Parallel global-route MPPI modes and long-horizon cooperative conflict prior & Prototype complete \\
V8 & Nav2 C++ controller integration through Phase~3 & In progress \\
\bottomrule
\end{tabular}

\section{Current conclusion}

The sandbox has moved from a single-gradient geodesic bias to an explicit
topology-informed multimodal proposal system. It can identify separate local
routes, translate them into dynamically constrained controls, preserve them as
separate MPPI sample groups, protect the global-path fallback, and apply a
tested Gaussian-mixture importance correction. The first genuine multimodal
V4 comparison exposed severe finite-sample degeneracy; the V4.1 ESS guard
restored success by reverting coherently to V3 when corrected evidence was
unreliable.

The PyTorch reciprocal pilot now additionally retains several global A* route
topologies as parallel, route-conditioned MPPI groups and evaluates eventual
peer conflict over the remaining announced routes. This is an implemented
prototype, not yet final evidence: route geometry is static within an episode,
and dynamic space--time topology generation remains open.

The project has now crossed from sandbox-only research into a working Nav2
integration. This establishes the implementation path, but it is not yet the
final thesis result. The central remaining question is whether topology-
conditioned ancillary proposals outperform a strong vanilla/path-biased Nav2
MPPI baseline in matched narrow and dynamic environments. Collision-validating
the ancillary rollouts, recording actual allocation/selection telemetry, and
running that benchmark are therefore higher priority than adding another
visual feature. Fixed covariance with V3-style grouped sampling remains the
trusted sandbox baseline, and raw V4 importance correction remains diagnostic
until broader experiments establish an advantage.

\end{document}
